\documentclass[11pt,a4paper]{article}
\usepackage[T1]{fontenc}
\usepackage[utf8]{inputenc}
\usepackage[lithuanian]{babel}
\usepackage{lmodern}
\usepackage{microtype}
\setlength{\emergencystretch}{1.5em}
\usepackage{amsmath,amssymb,array,tabularx,enumitem,fancyhdr}
\usepackage[a4paper,margin=1.7cm,headheight=15pt]{geometry}
\setlength{\parindent}{0pt}
\setlength{\parskip}{4pt}
\setlist[itemize]{leftmargin=1.5em,itemsep=1pt,topsep=3pt}
\pagestyle{fancy}\fancyhf{}
\lhead{Kvadratinės lygtys: vertinimo instrukcija}\rhead{Mokytojui}
\cfoot{\thepage}
\newcommand{\uzd}[3]{\subsection*{#1 uždavinys \hfill #2 taškai}\textit{Sąlyga.} #3\par}
\newcommand{\kriterijai}{\textbf{Taškų skyrimas:}}
\newcommand{\balas}[2]{\item \textbf{#1 t.} #2}
\begin{document}
\begin{center}
{\LARGE\bfseries Vertinimo instrukcija}\par
{\large Odisėjo variantas}\par
Varianto maksimumas -- \textbf{35 taškai}.
\end{center}
\textbf{Šaltinis.} Instrukcija parengta pagal naujausias Odisėjo varianto uždavinių sąlygas ir rekomenduojamą taškų paskirstymą. Lygiaverčiai teisingi sprendimo būdai vertinami tuo pačiu maksimumu.

\section*{Bendrosios taisyklės}
\begin{itemize}
\item Taškai skiriami už parodytą matematinį turinį, ne už formulių ar eilučių kiekį. Sujungti, bet aiškiai pagrįsti žingsniai gauna visus atitinkamus taškus.
\item Vertinama atskirai kiekviename \textbf{savarankiškame vertinimo vienete}: visas uždavinys, o uždaviniuose su dalimis -- kiekviena a), b), c) dalis. Taigi antroji klaida vienoje dalyje nepanaikina teisingai išspręstos kitos dalies taškų. Tai šioje instrukcijoje pasirinkta klaidų taisyklės taikymo apimtis.
\item \textbf{Viena aritmetinė klaida.} Jei po vienos skaičiavimo klaidos mokinys nuosekliai ir matematiškai teisingai tęsia sprendimą su savo gautais skaičiais, atimamas tik \textbf{1 taškas} tame vertinimo vienete. Tos klaidos sukeltas neteisingas galutinis atsakymas papildomai nebaudžiamas.
\item \textbf{Antroji nepriklausoma aritmetinė klaida.} Jei tame pačiame vienete pasitaiko dar viena neištaisyta aritmetinė klaida, skiriami tik taškai už teisingai atliktus kriterijus \textbf{iki pirmosios aritmetinės klaidos}. Už pirmąją klaidą dar vienas taškas papildomai neatimamas. Už vėlesnę sprendimo dalį taškų neskiriama.
\item \textbf{Klaidos pasekmė nėra nauja klaida.} Pavyzdžiui, klaidinga šaknų suma, iš jos teisingai apskaičiuota kvadratų suma ir iš jos teisingai sudaryta lygtis -- viena klaida. Nauja klaida būtų neteisingai suskaičiuota paties mokinio gauta skaičių išraiška.
\item Aiškiai ištaisyta klaida nelaikoma neištaisyta klaida. Vertinamas galutinis, aiškiai pasirinktas sprendimo kelias. Jei pateikti keli prieštaringi atsakymai ir nė vienas nepasirinktas, galutinio atsakymo kriterijus neįvykdytas.
\item \textbf{Metodinė klaida} (neteisinga formulė, neteisėtas dalijimas, neteisingas keitinys, nepagrįstas sprendinių atmetimas) nevertinama kaip vien aritmetinė. Taškai už klaidingą metodą ir nuo jo priklausomas nepagrįstas išvadas neskiriami; nepriklausomi teisingi žingsniai gali būti įskaityti.
\item Vienos klaidos taisyklė nesuteikia taškų už neparodytus gebėjimus: jei klaida uždavinį supaprastino ir mokinys nebetikrino šakų ar sąlygų, neparodyti kriterijai neįskaitomi. Vertinimo minimumas -- 0 taškų.
\end{itemize}

\textbf{Tik atsakymas.} Be sprendimo galima skirti tik aiškiai atsakymui numatytus taškus. Išimtis -- tiesioginis skaidymo dauginamaisiais užrašas: teisinga pilna sandauga pati parodo prašytą rezultatą ir gauna visus tai daliai skirtus taškus.



\newpage
\section*{Uždavinių patikra ir atsakymai}
Visi aštuoni uždaviniai matematiškai korektiški. Atsakymai yra sveikieji skaičiai, paprastosios trupmenos arba daugianarių sandaugos. 4 uždavinyje abu mokiniai tikrai pakeitė būtent nurodytą koeficientą. 5 uždavinyje gaunami keturi sveikieji sprendiniai. 6 uždavinio reiškinio reikšmė nepriklauso nuo pasirinktos trinario šaknies. 7 uždavinio sąlygos vienareikšmiškai nulemia dvi neigiamas šaknis.

\begin{center}\small\renewcommand{\arraystretch}{1.25}
\begin{tabularx}{\linewidth}{|c|X|}\hline
\textbf{Uždavinys}&\textbf{Atsakymas}\\\hline
1&$x=3$ arba $x=-\frac12$\\\hline
2&$x_1^3+x_2^3=45$\\\hline
3&a) $-(3x-4)(x-2)$; b) $(x+4)(3x-2y)$; c) $(x-3)(x^2+3x+9)$\\\hline
4&$b=-9$, $c=-10$\\\hline
5&$x\in\{-1,1,3,5\}$\\\hline
6&$11$\\\hline
7&$x_1=-5$, $x_2=-3$; $x^2+8x+15=0$\\\hline
8&$y^2-64y+324=0$\\\hline
\end{tabularx}\end{center}

\subsection*{Aritmetinių klaidų taisyklės pavyzdys: 8 uždavinys}
Teisingai reikia gauti $S=6$, $P=2$, naują sumą $64$ ir naują sandaugą $324$.
\begin{itemize}
\item Mokinys teisingai užrašo naujos sumos formulę, bet skaičiuodamas $36-4+24+8$ gauna $66$. Toliau teisingai apskaičiuoja sandaugą $324$ ir sudaro $y^2-66y+324=0$. \textbf{4/5 t.}: viena aritmetinė klaida, jos pasekmės papildomai nebaudžiamos.
\item Jei vėliau dar ir $18^2$ apskaičiuoja kaip $300$, įskaitomi tik kriterijai iki pirmosios klaidos: $S$ (1 t.) ir $P$ (1 t.). Vertinimas -- \textbf{2/5 t.}
\item Formulė $x_1^2+x_2^2=(x_1+x_2)^2$ yra metodinė klaida: trūksta $-2x_1x_2$ nario. Vieno taško aritmetinės klaidos lengvata čia netaikoma.
\end{itemize}

\textbf{Lygiaverčiai metodai.} Už teisingą trumpesnį kelią skiriamas visas balas. 2 ir 6 uždaviniuose nedraudžiama apskaičiuoti pradinių šaknų, todėl toks teisingas kelias gali būti vertinamas visu balu. 8 uždavinyje šaknų neskaičiavimo reikalavimas aiškiai nurodytas.

\textbf{Taškų suma.} $3+3+6+4+6+4+4+5=35$. Kiekvienas skaidymo uždavinio punktas vertinamas atskirai po 2 taškus.

\newpage
\section*{Odisėjo variantas: 1--3 uždaviniai}
\uzd{1}{3}{$(3x-2)^2=(x+4)^2$.}
\textbf{Sprendimas.} Kvadratų skirtumas duoda $[(3x-2)-(x+4)][(3x-2)+(x+4)]=0$, t. y. $(2x-6)(4x+2)=0$. Atsakymas: $x=3$ arba $x=-\frac12$.
\kriterijai\begin{itemize}
\balas{1}{Gauna teisingą kvadratinę lygtį arba kvadratų skirtumo sandaugą.}
\balas{1}{Išskaido arba gauna dvi teisingas tiesines lygtis.}
\balas{1}{Pateikia abu sprendinius $3,-1/2$.}
\end{itemize}
\textit{Vertinimo pastaba.} Galima spręsti $3x-2=x+4$ arba $3x-2=-(x+4)$. Būtina išnagrinėti abi šakas.
\uzd{2}{3}{$x^2-3x-2=0$; rasti šaknų kubų sumą.}
\textbf{Sprendimas.} $S=x_1+x_2=3$, $P=x_1x_2=-2$. Todėl $x_1^3+x_2^3=S^3-3PS=27+18=45$.
\kriterijai\begin{itemize}
\balas{1}{Pagal Vijeto teoremą gauna sumą $3$ ir sandaugą $-2$.}
\balas{1}{Taiko tapatybę $x_1^3+x_2^3=S^3-3PS$.}
\balas{1}{Apskaičiuoja rezultatą $45$.}
\end{itemize}
\textit{Vertinimo pastaba.} Tinka ir teisingas tiesioginis šaknų apskaičiavimas bei kubų sudėjimas.
\uzd{3}{6}{a) $-3x^2+10x-8$; b) $3x^2-2xy+12x-8y$; c) $x^3-27$.}
\textbf{Sprendimas.} \textbf{a)} $-3x^2+10x-8=-(3x-4)(x-2)$.\par\textbf{b)} $x(3x-2y)+4(3x-2y)=(x+4)(3x-2y)$.\par\textbf{c)} $x^3-3^3=(x-3)(x^2+3x+9)$.
\kriterijai\begin{itemize}
\balas{2}{a): 1 t. už tinkamas šaknis ar skaidymo pradžią; 1 t. už teisingą pilną sandaugą.}
\balas{2}{b): 1 t. už grupavimą su bendruoju daugikliu $3x-2y$; 1 t. už galutinę sandaugą.}
\balas{2}{c): 1 t. už kubų skirtumo formulės taikymą; 1 t. už teisingą sandaugą.}
\end{itemize}
\textit{Vertinimo pastaba.} Pilna teisinga sandauga iš karto gauna 2 t. atitinkamoje dalyje. Klaidų taisyklė taikoma kiekvienai daliai atskirai.
\newpage
\section*{Odisėjo variantas: 4--6 uždaviniai}
\uzd{4}{4}{Pirmasis neteisingai nurašė tik $c$ ir gavo $3,6$; antrasis neteisingai nurašė tik $b$ ir gavo $-2,5$. Abu savo lygtis išsprendė teisingai. Rasti $b,c$.}
\textbf{Sprendimas.} Pirmojo mokinio koeficientas $b$ teisingas: $-b=3+6=9$, todėl $b=-9$. Antrojo mokinio koeficientas $c$ teisingas: $c=(-2)\cdot5=-10$. Pradinė lygtis: $x^2-9x-10=0$.
\kriterijai\begin{itemize}
\balas{1}{Paaiškina, kad pirmojo mokinio šaknų suma nustato teisingą $b$.}
\balas{1}{Gauna $b=-9$.}
\balas{1}{Paaiškina, kad antrojo mokinio šaknų sandauga nustato teisingą $c$.}
\balas{1}{Gauna $c=-10$.}
\end{itemize}
\textit{Vertinimo pastaba.} Paaiškinimo taškams pakanka lygybių $-b=3+6$ ir $c=(-2)\cdot5$. Pirmasis sprendė $x^2-9x+18=0$, antrasis -- $x^2-3x-10=0$. Vien du teisingi atsakymai be pagrindimo -- iki 2/4 t.
\uzd{5}{6}{$2(2x^2-8x+3)(x^2-4x-3{,}5)=39$.}
\textbf{Sprendimas.} Imame $t=2x^2-8x$. Tuomet $2(x^2-4x-3{,}5)=t-7$. Gauname $(t+3)(t-7)=39$, $t^2-4t-60=0$, $(t-10)(t+6)=0$.\par Kai $t=10$, $x^2-4x-5=0$, $x=-1,5$. Kai $t=-6$, $x^2-4x+3=0$, $x=1,3$. Atsakymas: $x\in\{-1,1,3,5\}$.
\kriterijai\begin{itemize}
\balas{1}{Parenka keitinį $t=2x^2-8x$ arba lygiavertį $u=x^2-4x$.}
\balas{1}{Gauna lygtį $t^2-4t-60=0$ (arba $u^2-2u-15=0$).}
\balas{1}{Randa abi keitinio reikšmes $t=10,-6$ (arba $u=5,-3$).}
\balas{1}{Išsprendžia pirmą šaką ir gauna $x=-1,5$.}
\balas{1}{Išsprendžia antrą šaką ir gauna $x=1,3$.}
\balas{1}{Pateikia visus keturis sprendinius.}
\end{itemize}
\textit{Vertinimo pastaba.} Tinka ir $u=x^2-4x$: gaunama $u^2-2u-15=0$, $u=5,-3$. Praleidus vieną šaką, neįskaitomas jos ir pilno atsakymo kriterijus.
\uzd{6}{4}{$P(x)=x^2-4x-2$, $a$ -- viena šaknis; rasti $a^3-3a^2-6a+9$.}
\textbf{Sprendimas.} $P(a)=0$ duoda $a^2=4a+2$. Tada $a^3=4a^2+2a=18a+8$. Todėl\[a^3-3a^2-6a+9=(18a+8)-3(4a+2)-6a+9=11.\]\textbf{Trumpesnis kelias:} $a^3-3a^2-6a+9=(a+1)(a^2-4a-2)+11=(a+1)P(a)+11=11$.
\kriterijai\begin{itemize}
\balas{1}{Iš šaknies sąlygos gauna $P(a)=0$ arba $a^2=4a+2$.}
\balas{1}{Pradeda mažinti kubo laipsnį: $a^3=4a^2+2a$.}
\balas{1}{Gauna $a^3=18a+8$.}
\balas{1}{Teisingai įstato į reiškinį ir gauna reikšmę $11$.}
\end{itemize}
\textit{Vertinimo pastaba.} Teisingas trumpesnis kelias per $(a+1)P(a)+11$ gauna 4/4 t. Atsakymas vienodas abiem šaknims; atskirai jų tikrinti nereikia.
\newpage
\section*{Odisėjo variantas: 7--8 uždaviniai}
\uzd{7}{4}{$x_1<x_2<0$, $x_1^2+x_2^2=34$, $x_2-x_1=2$; sudaryti redukuotą lygtį su šaknimis $x_1,x_2$.}
\textbf{Sprendimas.} Pažymime $S=x_1+x_2$. Tuomet $S^2=2(x_1^2+x_2^2)-(x_2-x_1)^2=2\cdot34-2^2=64$. Kadangi abi šaknys neigiamos, $S=-8$. Sandauga $P=x_1x_2=\frac{S^2-(x_1^2+x_2^2)}{2}=\frac{64-34}{2}=15$. Ieškoma lygtis: $x^2+8x+15=0$, jos šaknys $-5$ ir $-3$.
\kriterijai\begin{itemize}
\balas{1}{Užrašo ryšį $S^2=2(x_1^2+x_2^2)-(x_2-x_1)^2$ arba lygiavertį.}
\balas{1}{Apskaičiuoja $S^2=64$ ir, atsižvelgęs į neigiamas šaknis, gauna $S=-8$.}
\balas{1}{Apskaičiuoja sandaugą $P=\frac{S^2-(x_1^2+x_2^2)}{2}=15$.}
\balas{1}{Sudaro lygtį $x^2+8x+15=0$.}
\end{itemize}
\textit{Vertinimo pastaba.} Taip pat tinka pirmiausia apskaičiuoti sandaugą iš $(x_2-x_1)^2=(x_1^2+x_2^2)-2x_1x_2$: gaunama $x_1x_2=15$, o tada iš sumos kvadrato nustatyti $x_1+x_2=-8$. Sąlyga $x_1<x_2<0$ būtina, kad būtų pasirinkta neigiama suma.
\uzd{8}{5}{$x^2-6x+2=0$ šaknys $x_1,x_2$; neskaičiuojant jų sudaryti redukuotą lygtį šaknims $(x_1+2)^2,(x_2+2)^2$.}
\textbf{Sprendimas.} $S=x_1+x_2=6$, $P=x_1x_2=2$. Nauja šaknų suma:\[(x_1+2)^2+(x_2+2)^2=S^2-2P+4S+8=36-4+24+8=64.\]Nauja sandauga:\[(x_1+2)^2(x_2+2)^2=(P+2S+4)^2=(2+12+4)^2=324.\]Atsakymas: $y^2-64y+324=0$.
\kriterijai\begin{itemize}
\balas{1}{Pagal Vijeto teoremą nustato $S=6$ ir $P=2$.}
\balas{1}{Išveda naujos sumos formulę $S^2-2P+4S+8$ arba lygiavertę formulę.}
\balas{1}{Gauna naują sumą $64$.}
\balas{1}{Išveda naujos sandaugos formulę $(P+2S+4)^2$ ir gauna $324$.}
\balas{1}{Užrašo redukuotą lygtį $y^2-64y+324=0$.}
\end{itemize}
\textit{Vertinimo pastaba.} Lygiavertis kelias: poslinkio šaknų suma $S+4=10$, sandauga $P+2S+4=18$, tada kvadratų suma $10^2-2\cdot18=64$, sandauga $18^2=324$. Kintamąjį galima žymėti ir $x$. Jei tik apskaičiuojamos pradinės šaknys ir jos naudojamos, nepaisant sąlygos, galima skirti iki 3 t. už naujos sumos ir sandaugos reikšmes bei lygtį, bet ne už prašytą šaknų neskaičiavimo metodą.
\subsection*{Pažymio lentelė}
\begin{center}\small\begin{tabular}{|c|*{10}{c|}}\hline
Taškai&0--3&4--7&8--10&11--14&15--17&18--21&22--24&25--28&29--31&32--35\\\hline
Pažymys&1&2&3&4&5&6&7&8&9&10\\\hline
\end{tabular}\end{center}
\end{document}